\label{sec:exp}

\subsection{Datasets and Evaluation Profiles}
We separate a gold-mask (``real'') profile used for primary claims from cross-dataset silver-label stress tests.
Table~\ref{tab:datasets} summarizes the profiles.

\input{tables/datasets_profiles}

\subsection{Baselines and Matched-Budget Protocol}
We compare tokenization strategies (e.g., fixed-grid and ROI-style variants) under a matched-cost protocol.
All comparisons are performed across a fixed set of budgets ($2\times 10^6$ to $7\times 10^6$ in six steps) with either FLOPs-matched or latency-aware matching depending on the claim.
We report quality and grounding together with latency and supportedness, and we enforce hard gates where appropriate (e.g., latency p95 and supportedness deltas).
To prevent post-hoc threshold tuning, calibration thresholds (e.g., $\tauRef$) are selected once on development data and frozen for test evaluation.

\subsection{Metrics}
Our evaluation emphasizes what is auditable:
\begin{itemize}
  \item \textbf{Grounding.} When gold masks are available (real profile), we report sentence-level grounding metrics (e.g., IoU) that connect citations to spatial support.
  \item \textbf{Supportedness.} We report an \emph{unsupported rate} defined as the fraction of frames flagged by the verifier as unsupported (issue U1).
  \item \textbf{Clinical correctness proxies.} We report frame-level correctness on structured findings (frame F1). For safety-critical findings, we additionally report \emph{critical-present recall} averaged over studies with at least one ground-truth critical present frame (the critical set is fixed and audited).
  \item \textbf{Refusal.} We report refusal rate and refusal calibration (ECE/reliability), and we enforce an anti-silencing constraint via \emph{critical miss-rate} (ground-truth critical findings that are refused).
\end{itemize}

\subsection{Statistical Protocol}
We follow paired bootstrap testing (typically $n_{\mathrm{boot}}{=}20{,}000$) with a one-sided alternative when comparing a method against a baseline, and we apply Holm correction over the relevant family (e.g., multiple budgets or counterfactual variants) to control family-wise error.
We report confidence intervals and explicitly track the number of random seeds used for each artifact.
Implementation details and artifact paths are provided in Table~\ref{tab:oral-minset} and Appendix tables.

To avoid ``multiple comparisons by choice,'' we pre-register hypothesis families at the \emph{claim level}.
Each claim fixes: the evaluation profile (real vs. silver stress test), the budget sweep or perturbation set, the primary metric(s), the direction of the alternative, and the multiple-comparison correction scope.
Exact families and pass rules are summarized in Appendix Table~\ref{tab:stats-family} and implemented in our audit scripts.

Concretely, we use:
\begin{itemize}
  \item \textbf{Budget families.} For multi-budget claims (e.g., C0001/C0004), we treat the budget sweep as the family for each pre-registered hypothesis (per metric; and per baseline when multiple baselines are required) and apply Holm correction over budgets.
  \item \textbf{Counterfactual families.} For counterfactual non-triviality (e.g., C0003 and V0003), we treat perturbation variants (e.g., no\_cite, cite\_swap, $\omega$-perm) as the family within each metric group and apply Holm correction accordingly.
\end{itemize}

We treat different claims as different hypothesis families because they address different questions (tradeoffs, grounding, counterfactual non-triviality, refusal calibration), but we use the same fixed splits and do not tune thresholds on the test set.

\subsection{External Validity}
\paragraph{Backbone transfer.}
To test whether improvements are architecture-specific, we replicate key evaluations under multiple report-generation backbones (e.g., a lightweight generator and an LLM-backed PCG), while keeping the matched-budget protocol and refusal calibration fixed.
In particular, refusal thresholds and verifier rules are frozen to prevent post-hoc tuning on each backbone.

\paragraph{Dataset shift.}
We separate primary gold-mask claims (real profile) from cross-dataset stress tests that rely on silver labels (automatic or pseudo masks).
Under dataset shift, we use counterfactual tests (e.g., $\omega$-permutation) to verify that the citation channel remains non-trivial under guarded statistical protocols.
